\begin{question}{30}{
    De Koninklijke Marine vergelijkt twee sonarsystemen, $A$ en $B$, op haar fregatten.
    Tegen hetzelfde type doel is de detectieafstand (in \si{\kilo\metre}) gemeten bij acht oefeningen met systeem $A$ en tien oefeningen met systeem $B$.
    Neem aan dat de afstanden normaal verdeeld zijn met onbekende gemiddelden $\mu_A$ en $\mu_B$ en standaardafwijkingen $\sigma_A$ en $\sigma_B$.

    \begin{center}
        \begin{tabular}{m{0.24\textwidth}|cccccccccccc}
            \toprule
                \textbf{Detectieafstanden $A$} & $18.2$ & $17.5$ & $18.8$ & $17.9$ & $18.4$ & $17.7$ & $18.1$ & $18.6$ &        &        \\
            \midrule
                \textbf{Detectieafstanden $B$} & $16.5$ & $19.4$ & $15.2$ & $18.8$ & $14.9$ & $17.6$ & $19.9$ & $16.1$ & $18.3$ & $15.8$ \\
            \bottomrule
        \end{tabular}
    \end{center}
    Je mag in deze vraag bij hypothesetoetsen gebruik maken van een significantieniveau $\alpha = 0.05$.
}

    \subquestion{6}{
        Bereken de steekproefgemiddelden en de steekproefvarianties van de detectieafstanden van beide sonarsystemen.
    }
    \solution{
    
        We berekenen het steekproefgemiddelde $\overline{x_A}$ voor de $n = 8$ detectieafstanden van sonarsysteem $A$ als volgt:
        \begin{align*}
            \overline{x_A} = \frac{a_1+a_2+\ldots+a_n}{n} = \frac{ 18.2 + 17.5 + \ldots + 18.6 }{ 8 } = 18.15.\rubric{1}
        \end{align*}
    
        We berekenen de steekproefvariantie $s_A^2$ als volgt:
        \begin{align*}
            s_A^2 &= \frac{ (a_1 - \overline{x_A})^2 + (a_2 - \overline{x_A})^2 + \ldots + (a_n - \overline{x_A})^2 }{ n - 1 } \\
                &= \frac{ (18.2 - 18.15)^2 + (17.5 - 18.15)^2 + \ldots + (18.6 - 18.15)^2 }{ 8 - 1 } \\
                &\approx 0.1971.\rubric{3}
        \end{align*}

        We berekenen het steekproefgemiddelde $\overline{x_B}$ voor de $m = 10$ detectieafstanden van sonarsysteem $A$ als volgt:
        \begin{align*}
            \overline{x_B} = \frac{b_1+b_2+\ldots+b_m}{m} = \frac{ 16.5 + 19.4 + \ldots + 15.8 }{ 10 } = 17.25.\rubric{1}
        \end{align*}
    

        We berekenen de steekproefvariantie $s_B^2$ als volgt:
        \begin{align*}
            s_B^2 &= \frac{ (b_1 - \overline{x_B})^2 + (b_2 - \overline{x_B})^2 + \ldots + (b_m - \overline{x_B})^2 }{ m - 1 } \\
                &= \frac{ (16.5 - 17.25)^2 + (19.4 - 17.25)^2 + \ldots + (15.8 - 17.25)^2 }{ 10 - 1 } \\ \rubric{2}
                &\approx 3.2206.
        \end{align*}
    }

    \subquestion{10}{
        Voer een $F$-toets uit om te bepalen of de varianties in de detectieafstanden gelijk kunnen zijn voor beide sonarsystemen.
        Geef hierbij expliciet de stappen aan van het algemene stappenplan voor hypothesetoetsen.
        Bepaal de toetsuitslag op basis van het kritieke gebied.
    }
    \solution{
        \subsubsection*{Stap 1: formuleer de nulhypothese $H_0$ en de alternatieve hypothese $H_1$.}
        Bij een $F$-toets toetsen we of de varianties $\sigma_A^2$ en $\sigma_B^2$ gelijk zouden kunnen zijn, oftewel
        \begin{align*}
            H_0: \quad \sigma_A^2 = \sigma_B^2 \\
            H_1: \quad \sigma_A^2 \neq \sigma_B^2 \rubric{1}
        \end{align*}

        \subsubsection*{Stap 2: bepaal het significantieniveau $\alpha$.}
        Er is in de vraag gegeven dat we mogen werken met een significantieniveau $\alpha=0.05$.
        
        \subsubsection*{Stap 3: verzamel data voor het berekenen van de toetsingsgrootheid.}
        Ook zijn er in de vraag al steekproefdata verzameld voor beide populaties. 
        Op basis van deze steekproefdata hebben we in vraag (a) bepaald dat:
        
        \begin{center}
            \begin{tabular}{C{.3\textwidth} C{.3\textwidth} C{.3\textwidth}}
                $\overline{x_A} = 18.15$ & $s_A^2 \approx 0.1971$ & $n = 8$ \\
                $\overline{x_B} = 17.25$ & $s_B^2 \approx 3.2206$ & $m = 10$
            \end{tabular}
        \end{center}

        \subsubsection*{Stap 4: bereken de toetsingsgrootheid.}
        De theoretische toetsingsgrootheid (hoofdletter!) bij een $F$-toets is gelijk aan
        \[
            F = \frac{S_A^2}{S_B^2}, \rubric{1}
        \]
        en volgt onder $H_0$ een $F$-verdeling met $\dfn = n-1 = 7$ en $\dfd = m-1 = 9$ vrijheidsgraden. \rubric{1}

        De geobserveerde toetsingsgrootheid (kleine letter!) is gelijk aan
        \[
            f = \frac{s_A^2}{s_B^2} = \frac{0.1971}{3.2206} \approx 0.0612. \rubric{1}
        \]
        
        \subsubsection*{Stap 5: bepaal de toetsuitslag (met behulp van het kritieke gebied).}

        Een $F$-toets voor gelijke varianties is altijd tweezijdig, oftewel het kritieke gebied is van de vorm $(0, g_1)$ en $(g_2, \infty)$, waarbij de grenswaarden $g_1$ en $g_2$ als volgt kunnen worden berekend:
        \begin{enumerate}
            \item Bepaal de grenswaarde $g_1$ waarvoor geldt dat $P(F \leq g_1) = \alpha/2 = 0.025$.
            Voer hiervoor in de \enquote{numerieke solver} van de grafische rekenmachine in:
            \begin{align*}
                E_1 &= \fcdf(\lb=0, \ub=X, \df_1=7, \df_2=9) \\
                E_2 &= 0.025
            \end{align*}
            De \enquote{numerieke solver} geeft een waarde van $g_1 \approx 0.2073$.
            \item Bepaal de grenswaarde $g_2$ waarvoor geldt dat $P(F \geq g_2) = \alpha/2 = 0.025$.
            Voer hiervoor in de \enquote{numerieke solver} van de grafische rekenmachine in:
            \begin{align*}
                E_1 &= \fcdf(\lb=X, \ub=\GRinfinity, \df_1=7, \df_2=9) \\
                E_2 &= 0.025
            \end{align*}
            De \enquote{numerieke solver} geeft een waarde van $g_2 \approx 4.1970$.
        \end{enumerate}
        
        Het kritieke gebied is dus gelijk aan $(0, 0.2073)$ en $(4.1970, \infty)$. \rubric{4}

        De geobserveerde toetsingsgrootheid $f \approx 0.0612$ ligt in het kritieke gebied, dus we besluiten om $H_0$ te verwerpen.\rubric{1}

        \subsubsection*{Stap 6: vertaal de toetsuitslag terug naar de originele context van het probleem.}
        Er is voldoende reden om op basis van deze steekproeven de aanname van gelijke varianties (van de detectieafstanden van beide sonarsystemen) te verwerpen, 
        oftewel we gaan er nu van uit dat $\sigma_X^2 \neq \sigma_Y^2$. \rubric{1}
    }

    \subquestion{10}{
        Bepaal met behulp van een geschikte $t$-toets of de gemiddelde detectieafstand significant groter is voor sonarsysteem $A$.
        Geef hierbij expliciet de stappen aan van het algemene stappenplan voor hypothesetoetsen.
        Bepaal hierbij de toetsuitslag met behulp van de $p$-waarde.
    }
    \solution{
        \subsubsection*{Stap 1: formuleer de nulhypothese $H_0$ en de alternatieve hypothese $H_1$.}
        We gaan toetsen of de gemiddelde detectieafstanden $\mu_A$ significant groter is dan $\mu_B$.
        De hypothesen zijn daarom als volgt:
        \begin{align*}
            H_0: &\quad \mu_A \leq  \mu_B \text{ (niet significant grotere gemiddelde voor $A$) } \\
            H_1: &\quad \mu_A > \mu_B \text{ (wel significant grotere gemiddelde voor $A$) } \rubric{2}
        \end{align*}

        \subsubsection*{Stap 2: bepaal het significantieniveau $\alpha$.}
        Er is opnieuw gegeven dat we mogen werken met een significantieniveau $\alpha=0.05$.
        
        \subsubsection*{Stap 3: verzamel data voor het berekenen van de toetsingsgrootheid}
        Ook zijn er in de vraag al steekproefdata verzameld voor beide populaties. 
        Op basis van deze steekproefdata hebben we in vraag (a) bepaald dat:
        
        \begin{center}
            \begin{tabular}{C{.3\textwidth} C{.3\textwidth} C{.3\textwidth}}
                $\overline{x_A} = 18.15$ & $s_A^2 \approx 0.1971$ & $n = 8$ \\
                $\overline{x_B} = 17.25$ & $s_B^2 \approx 3.2206$ & $m = 10$
            \end{tabular}
        \end{center}

        \subsubsection*{Stap 4: bereken de toetsingsgrootheid.}
        Op basis van ons antwoord bij vraag (b) mogen we aannemen dat de varianties ongelijk zijn, oftewel $\sigma_X^2 \neq \sigma_Y^2$.
        De theoretische toetsingsgrootheid (hoofdletter!) bij een $t$-toets met twee steekproeven (met ongelijke varianties) is gelijk aan
        \[
            T = \frac{(\overline{X_A} - \overline{X_B}) - (\mu_A - \mu_B)}{\sqrt{\frac{S_A^2}{n} + \frac{S_B^2}{m}}}, \rubric{1}
        \]
        en volgt onder $H_0$ een $t$-verdeling met $\df=\min(n-1,m-1)=7$ vrijheidsgraden. \rubric{1}

        De geobserveerde toetsingsgrootheid (kleine letter!) is gelijk aan
        \begin{align*}
            t   &= \frac{(\overline{x_A}-\overline{x_B}) - (\mu_A - \mu_B)}{\sqrt{\frac{s_A^2}{n} + \frac{s_B^2}{m}}} \\
                &= \frac{(18.15 - 17.25) - 0}{\sqrt{\frac{0.1971}{8} + \frac{3.2206}{10}}} \\
                &\approx 1.5285. \rubric{2}
        \end{align*}

        \subsubsection*{Stap 5: bepaal de toetsuitslag (met behulp van de $p$-waarde).}
        Omdat we rechtszijdig toetsen, is de $p$-waarde gelijk aan $P(T \geq t)$:
        Deze kans kunnen we met de $t$-verdeling bepalen, namelijk:
        \begin{align*}
            P(T \geq t) &= \tcdf(\lb = 1.5285, \ub=\GRinfinity, \df=7) \approx 0.0851.
        \end{align*}
        De $p$-waarde is dus gelijk aan $p \approx 0.0851$. \rubric{2}

        Aangezien de $p$-waarde strikt groter is dan $\alpha$, besluiten we om de nulhypothese niet te verwerpen.\rubric{1}
        
        \subsubsection*{Stap 6: vertaal de toetsuitslag terug naar de originele context van het probleem.}
        Er is op basis van deze steekproeven onvoldoende bewijs om aan te nemen dat de gemiddelde detectieafstand bij sonarsysteem $A$
        significant groter is dan bij sonarsysteem $B$. \rubric{1}    
        % \begin{center}
        %     \resizebox{0.9\textwidth}{!}{
        %         \includegraphics{exam20260724-two-sample-t-test-pvalue.png} 
        %     }
        % \end{center}
    }

    \subquestion{4}{
        Is er op basis van de twee toetsen een voorkeur voor één van de systemen aan te wijzen?
        Licht toe aan de hand van de gegeven steekproefdata en de conclusies van de twee toetsen.
    }
    \solution{
        Uit de $t$-toets is gebleken dat we niet kunnen zeggen dat de gemiddelden significant van elkaar verschillen.\rubric{1}

        Uit de $F$-toets echter bleek wel dat de spreiding significant verschilt; in het bijzonder geven de steekproefdata 
        aan dat de spreiding bij $B$ significant groter dan bij $A$.\rubric{1}

        Systeem $A$ is daardoor meer \enquote{voorspelbaar} qua performance, en $B$ levert soms zeer lage detectieafstanden (\SI{14.9}{\kilo\metre}),
        wat operationeel riskant is.\rubric{1}
        
        De voorkeur gaat dus uit naar $A$.\rubric{1}
    }
\end{question}